\section{Tabel Kebenaran}

Tabel kebenaran (truth table) mendaftar secara sistematis semua kombinasi nilai kebenaran variabel proposisional dan nilai yang dihasilkan oleh proposisi majemuk \cite{rosen2019,forallx}. Matematika biasa menggunakan logika dua-nilai: setiap pernyataan bernilai benar atau salah. Tabel kebenaran memungkinkan kita melihat bagaimana kebenaran atau kesalahan pernyataan majemuk bergantung pada kebenaran atau kesalahan komponen-komponennya \cite{brilliant_truthtables}.

Untuk $n$ variabel proposisional, tabel kebenaran memerlukan $2^n$ baris agar semua kombinasi tercakup. Sebagai contoh, dua variabel ($p$, $q$) memerlukan 4 baris; tiga variabel memerlukan 8 baris. Urutan baris umumnya disusun secara sistematis (misalnya, dengan menganggap B=1 dan S=0, urutan biner 00, 01, 10, 11) untuk memastikan tidak ada kombinasi yang terlewat \cite{levin_discrete,brilliant_truthtables}. Grafik banyak baris $2^n$ untuk $n$ variabel dapat dilihat pada Gambar~\ref{fig:baris-tabel-kebenaran}.

\begin{figure}[htbp]
\centering
\includegraphics[width=0.75\textwidth]{baris-tabel-kebenaran}
\caption{Banyak baris tabel kebenaran: $2^n$ untuk $n$ variabel proposisional (dihasilkan oleh script Python)}
\label{fig:baris-tabel-kebenaran}
\end{figure}

Grafik di atas menggambarkan pertumbuhan eksponensial jumlah baris ketika variabel bertambah.

\begin{contoh}
Tabel kebenaran untuk operator dasar $\neg p$, $p \wedge q$, $p \vee q$, $p \rightarrow q$, $p \leftrightarrow q$:

\begin{center}
\small
\begin{tabular}{|c|c|c|c|c|c|c|}
\hline
$p$ & $q$ & $\neg p$ & $p \wedge q$ & $p \vee q$ & $p \rightarrow q$ & $p \leftrightarrow q$ \\
\hline
B & B & S & B & B & B & B \\
B & S & S & S & B & S & S \\
S & B & B & S & B & B & S \\
S & S & B & S & S & B & B \\
\hline
\end{tabular}
\end{center}
\end{contoh}

Tabel kebenaran untuk implikasi $p \rightarrow q$ kerap menimbulkan pertanyaan. Baris ketiga dan keempat (ketika $p$ salah) menghasilkan nilai benar. Alasannya, implikasi material menangkap gagasan: ``tidak pernah terjadi $p$ benar dan $q$ salah.'' Jika $p$ salah, syarat tersebut tidak dilanggar, sehingga implikasi dianggap benar \cite{brilliant_truthtables,libretexts_logic}.

\begin{contoh}
Tabel kebenaran untuk $(p \rightarrow q) \wedge (q \rightarrow p)$ dan $p \leftrightarrow q$ (keduanya ekuivalen):

\begin{center}
\small
\begin{tabular}{|c|c|c|c|c|c|}
\hline
$p$ & $q$ & $p \rightarrow q$ & $q \rightarrow p$ & $(p \rightarrow q) \wedge (q \rightarrow p)$ & $p \leftrightarrow q$ \\
\hline
B & B & B & B & B & B \\
B & S & S & B & S & S \\
S & B & B & S & S & S \\
S & S & B & B & B & B \\
\hline
\end{tabular}
\end{center}
Kolon ke-5 dan ke-6 identik: bikondisional $p \leftrightarrow q$ memang ekuivalen dengan konjungsi dua implikasi.
\end{contoh}

Metode tabel kebenaran dikaitkan dengan Ludwig Wittgenstein dan Emil Post, meskipun ide-ide pendahulunya telah ada pada Frege, Russell, Boole, dan Peirce \cite{wiki_proplogic}. Alat ini sangat berguna untuk memeriksa kevalidan argumen, menentukan apakah suatu formula adalah tautologi (selalu benar) atau kontradiksi (selalu salah), dan untuk merancang serta menganalisis sirkuit digital \cite{uw_boolean,levin_discrete}.

Dalam ilmu komputer, tabel kebenaran berkaitan erat dengan gerbang logika (AND, OR, NOT, dan sebagainya). Setiap gerbang memiliki tabel kebenaran yang mendefinisikan keluarannya untuk setiap kombinasi masukan. Pengetahuan tentang tabel kebenaran membantu dalam perancangan rangkaian kombinatorial dan pemahaman perilaku program yang bergantung pada kondisi Boolean \cite{grinnell_logic,brilliant_truthtables}.

\begin{contoh}
\textbf{Generator tabel kebenaran dengan Python} \cite{python_docs_boolean,levin_discrete}

Fungsi berikut menghasilkan semua kombinasi nilai kebenaran untuk $n$ variabel dan mencetak tabel untuk operator implikasi $p \to q$:

\begin{lstlisting}[style=pythonstyle, caption={Kombinasi nilai kebenaran dan tabel untuk $p \to q$}]
from itertools import product

def truth_table_impl(n=2):
    """Cetak kombinasi B/S untuk n variabel dan kolom p -> q."""
    vars_list = [f"p{i}" for i in range(1, n + 1)]
    print("  ".join(vars_list), "  p -> q")
    print("-" * (4 * n + 8))
    for row in product([True, False], repeat=n):
        p, q = row[0], row[1] if n > 1 else False
        impl = (not p) or q  # p -> q
        row_str = "  ".join("B" if x else "S" for x in row)
        print(row_str, "  ", "B" if impl else "S")

truth_table_impl(2)  # 4 baris untuk p, q
\end{lstlisting}
\end{contoh}
